\documentclass[12pt,a4paper]{article}

% Paquetes esenciales
\usepackage[utf8]{inputenc}
\usepackage[spanish]{babel}
\usepackage[T1]{fontenc}
\usepackage{lmodern}

% Geometría de la página
\usepackage[left=2.5cm,right=2.5cm,top=2.5cm,bottom=2.5cm]{geometry}

% Matemáticas
\usepackage{amsmath}
\usepackage{amssymb}
\usepackage{amsfonts}

% Tablas
\usepackage{booktabs}
\usepackage{longtable}
\usepackage{array}
\usepackage{multirow}
\usepackage{tabularx}

% Gráficos y colores
\usepackage{graphicx}
\usepackage{xcolor}
\usepackage{colortbl}

% Código fuente
\usepackage{listings}
\usepackage{courier}

% Enlaces e hipervínculos
\usepackage{hyperref}
\hypersetup{
    colorlinks=true,
    linkcolor=blue,
    filecolor=magenta,      
    urlcolor=cyan,
    citecolor=blue,
    pdftitle={Bitácora del Proceso de Cifrado de Imágenes},
    pdfauthor={Sistema de Documentación Automática},
    pdfsubject={Cifrado Caótico de Imágenes},
    pdfkeywords={Cifrado, Caos, Rössler, Mapa Logístico, MQTT, TLS}
}

% Otros paquetes útiles
\usepackage{fancyhdr}
\usepackage{titlesec}
\usepackage{enumitem}
\usepackage{float}

% Configuración de listings para Python
\lstset{
    language=Python,
    basicstyle=\ttfamily\small,
    keywordstyle=\color{blue}\bfseries,
    commentstyle=\color{green!60!black}\itshape,
    stringstyle=\color{red},
    numbers=left,
    numberstyle=\tiny\color{gray},
    stepnumber=1,
    numbersep=8pt,
    backgroundcolor=\color{gray!10},
    showspaces=false,
    showstringspaces=false,
    showtabs=false,
    frame=single,
    tabsize=4,
    captionpos=b,
    breaklines=true,
    breakatwhitespace=false,
    escapeinside={(*@}{@*)},
    xleftmargin=2em,
    framexleftmargin=1.5em
}

% Configuración de encabezados y pies de página
\pagestyle{fancy}
\fancyhf{}
\fancyhead[L]{\leftmark}
\fancyhead[R]{\thepage}
\renewcommand{\headrulewidth}{0.4pt}

% Profundidad del índice
\setcounter{tocdepth}{3}
\setcounter{secnumdepth}{3}

% Información del documento
\title{\textbf{Bitácora del Proceso de Cifrado de Imágenes}\\
\large{Código: Maestro\_TLS\_Keystream\_XYZ.py}}
\author{Sistema de Documentación Automática}
\date{13 de febrero de 2026 \\ Versión 1.0}

\begin{document}

\maketitle
\thispagestyle{empty}

\newpage
\tableofcontents
\newpage

\section{Resumen Ejecutivo}

Este documento describe detalladamente el proceso de cifrado de imágenes implementado en \texttt{Maestro\_TLS\_Keystream\_XYZ.py}, el cual se ejecuta en una Raspberry Pi 4. El sistema utiliza técnicas criptográficas avanzadas basadas en \textbf{caos determinístico} para cifrar imágenes RGB mediante dos etapas principales:

\begin{enumerate}
    \item \textbf{Difusión} (usando el Mapa Logístico)
    \item \textbf{Confusión} (usando el Oscilador de Rössler)
\end{enumerate}

Una vez cifrada la imagen, los datos se transmiten de forma segura a una segunda Raspberry Pi utilizando el protocolo \textbf{MQTT con capa TLS}.

\subsection*{Contexto de la imagen de ejemplo:}
\begin{itemize}
    \item Tipo: RGB (3 canales)
    \item Dimensiones: $250 \times 250$ píxeles
    \item Total de píxeles: 62,500
    \item Total de valores: 187,500 ($250 \times 250 \times 3$)
\end{itemize}

\section{FASE DE INICIALIZACIÓN Y CONFIGURACIÓN}

\subsection{Parámetros del Sistema}

El sistema se configura con los siguientes parámetros críticos:

\subsubsection{Parámetros del Oscilador de Rössler}

\begin{lstlisting}
a = 0.2
b = 0.2
c = 5.7
\end{lstlisting}

Estos parámetros controlan el comportamiento caótico del sistema de Rössler, garantizando que las trayectorias generadas sean impredecibles y sensibles a las condiciones iniciales.

\subsubsection{Condiciones Iniciales de Rössler}

\begin{lstlisting}
Y0 = [0.1, 0.1, 0.1]  # [x0, y0, z0]
\end{lstlisting}

\subsubsection{Parámetros del Mapa Logístico}

\begin{lstlisting}
aLog = 3.99  # Parametro de control (debe estar en rango caotico: 3.57 < a <= 4)
x0_log = 0.4  # Condicion inicial (0 < x0 < 1)
\end{lstlisting}

\subsubsection{Parámetros de Integración y Sincronización}

\begin{lstlisting}
TIEMPO_SINC = 6000   # Iteraciones de sincronizacion
KEYSTREAM = 30000     # Iteraciones para generacion de clave
H = 0.01              # Paso de integracion
IMG_SCALE = 0.01      # Factor de escala para difusion
\end{lstlisting}

\subsubsection{Configuración MQTT con TLS}

\begin{lstlisting}
BROKER = "raspberrypiJED.local"
PORT = 8883           # Puerto seguro MQTT sobre TLS
QOS = 1               # Calidad de servicio
USERNAME = "usuario1"
PASSWORD = "qwerty123"
CA_CERT_PATH = "/etc/mosquitto/ca_certificates/ca.crt"
\end{lstlisting}

\section{FASE DE CARGA Y VECTORIZACIÓN DE LA IMAGEN}

\subsection{Proceso de Carga}

La función \texttt{cargar\_imagen()} realiza los siguientes pasos:

\begin{enumerate}
    \item \textbf{Lectura del archivo}: Se carga la imagen \texttt{Prueba2.jpg} usando PIL (Python Imaging Library)
    \item \textbf{Conversión a array NumPy}: La imagen se convierte en una matriz 3D de valores enteros (0-255)
    \item \textbf{Extracción de dimensiones}: Se obtienen alto, ancho y número de canales
    \item \textbf{Aplanamiento}: La matriz 3D se convierte en un vector 1D
    \item \textbf{Normalización}: Los valores se dividen entre 255.0 para obtener el rango [0, 1]
\end{enumerate}

\subsection{Representación Matemática}

Para una imagen RGB de $250 \times 250$:

\begin{equation}
\begin{aligned}
\text{Imagen original (3D)} &: [250, 250, 3] \\
\text{Vector aplanado (1D)} &: [187500] \\
\text{Rango de valores} &: [0.0, 1.0]
\end{aligned}
\end{equation}

\subsection{Ejemplo de Vectorización}

\begin{table}[H]
\centering
\caption{Ejemplo de los primeros 10 valores del vector normalizado}
\label{tab:vectorizacion}
\begin{tabular}{@{}ccccc@{}}
\toprule
\textbf{Índice} & \textbf{Posición Pixel} & \textbf{Canal} & \textbf{Valor Original (0-255)} & \textbf{Valor Normalizado (0-1)} \\ \midrule
0 & (0,0) & R & 145 & 0.568627 \\
1 & (0,0) & G & 87 & 0.341176 \\
2 & (0,0) & B & 203 & 0.796078 \\
3 & (0,1) & R & 156 & 0.611765 \\
4 & (0,1) & G & 92 & 0.360784 \\
5 & (0,1) & B & 198 & 0.776471 \\
6 & (0,2) & R & 134 & 0.525490 \\
7 & (0,2) & G & 76 & 0.298039 \\
8 & (0,2) & B & 211 & 0.827451 \\
9 & (0,3) & R & 167 & 0.654902 \\ \bottomrule
\end{tabular}
\end{table}

\textbf{Interpretación}: Cada píxel RGB se descompone en 3 valores consecutivos en el vector (R, G, B). La normalización facilita las operaciones matemáticas posteriores.

\section{FASE DE DIFUSIÓN (Mapa Logístico)}

\subsection{Objetivo de la Difusión}

La \textbf{difusión} es una técnica criptográfica que tiene como objetivo \textbf{dispersar la información} de la imagen original de manera que pequeños cambios en la entrada produzcan grandes cambios en la salida. Esto se logra mediante la \textbf{permutación} de los píxeles según una secuencia caótica.

\subsection{Generación del Mapa Logístico}

El mapa logístico es un sistema dinámico discreto definido por:

\subsubsection*{Ecuación del Mapa Logístico:}
\begin{equation}
x_{n+1} = a \cdot x_n \cdot (1 - x_n)
\label{eq:logistic}
\end{equation}

Donde:
\begin{itemize}
    \item $a = 3.99$ (parámetro de control en régimen caótico)
    \item $x_0 = 0.4$ (condición inicial)
    \item $n$ es el número de iteración
\end{itemize}

\subsubsection*{Propiedades clave:}
\begin{itemize}
    \item Genera una secuencia pseudoaleatoria determinística
    \item Extremadamente sensible a condiciones iniciales
    \item Los valores se mantienen en el rango $(0, 1)$
\end{itemize}

\subsection{Ejemplo de Secuencia Logística}

\begin{table}[H]
\centering
\caption{Primeros 10 valores del vector logístico}
\label{tab:logistic}
\begin{tabular}{@{}cll@{}}
\toprule
\textbf{Iteración} & \textbf{Valor $x_n$} & \textbf{Cálculo} \\ \midrule
0 & 0.400000 & (condición inicial) \\
1 & 0.956400 & $3.99 \times 0.4 \times (1 - 0.4)$ \\
2 & 0.166641 & $3.99 \times 0.9564 \times (1 - 0.9564)$ \\
3 & 0.554380 & $3.99 \times 0.166641 \times (1 - 0.166641)$ \\
4 & 0.986604 & $3.99 \times 0.55438 \times (1 - 0.55438)$ \\
5 & 0.052721 & $3.99 \times 0.986604 \times (1 - 0.986604)$ \\
6 & 0.199491 & $3.99 \times 0.052721 \times (1 - 0.052721)$ \\
7 & 0.637320 & $3.99 \times 0.199491 \times (1 - 0.199491)$ \\
8 & 0.922821 & $3.99 \times 0.63732 \times (1 - 0.63732)$ \\
9 & 0.284316 & $3.99 \times 0.922821 \times (1 - 0.922821)$ \\ \bottomrule
\end{tabular}
\end{table}

\subsection{Generación del Vector de Mezcla}

El vector logístico se transforma en índices enteros para crear el vector de mezcla:

\subsubsection*{Transformación:}
\begin{equation}
\text{vector\_mezcla}[i] = \lfloor \text{vector\_logistico}[i] \times n_{\text{max}} \rfloor
\label{eq:mixing}
\end{equation}

Donde $n_{\text{max}} = 187500$ (total de valores en la imagen $250 \times 250 \times 3$)

\begin{table}[H]
\centering
\caption{Conversión a índices de mezcla (primeros 10 valores)}
\label{tab:mixing}
\begin{tabular}{@{}cccc@{}}
\toprule
\textbf{Índice} & \textbf{Valor Logístico} & \textbf{Cálculo ($\times$ 187500)} & \textbf{Vector Mezcla (índice)} \\ \midrule
0 & 0.400000 & $0.4 \times 187500$ & 75000 \\
1 & 0.956400 & $0.9564 \times 187500$ & 179325 \\
2 & 0.166641 & $0.166641 \times 187500$ & 31245 \\
3 & 0.554380 & $0.55438 \times 187500$ & 103946 \\
4 & 0.986604 & $0.986604 \times 187500$ & 184988 \\
5 & 0.052721 & $0.052721 \times 187500$ & 9885 \\
6 & 0.199491 & $0.199491 \times 187500$ & 37404 \\
7 & 0.637320 & $0.63732 \times 187500$ & 119497 \\
8 & 0.922821 & $0.922821 \times 187500$ & 173029 \\
9 & 0.284316 & $0.284316 \times 187500$ & 53309 \\ \bottomrule
\end{tabular}
\end{table}

\textbf{Interpretación}: Cada valor del vector de mezcla representa un índice en el vector original de la imagen. Estos índices determinan el orden en que se leerán los píxeles.

\subsection{Proceso de Permutación (Algoritmo de Difusión)}

El algoritmo de difusión utiliza un \textbf{marcador especial (260.0)} para llevar un control de qué valores ya han sido utilizados. Se realizan dos pasadas:

\subsubsection{Primera Pasada: Lectura Aleatoria}
\begin{lstlisting}
for i in range(nmax):
    pos = vector_mezcla[i]
    if vector_temp[pos] != 260.0:
        difusion[contador] = vector_temp[pos]
        contador += 1
        vector_temp[pos] = 260.0  # Marcamos como usado
\end{lstlisting}

\subsubsection{Segunda Pasada: Lectura Secuencial de Restantes}
\begin{lstlisting}
for j in range(nmax):
    if contador >= nmax:
        break
    if vector_temp[j] != 260.0:
        difusion[contador] = vector_temp[j]
        contador += 1
\end{lstlisting}

\subsection{Ejemplo de Permutación}

\begin{table}[H]
\centering
\caption{Ejemplo del proceso de difusión (primeros 10 valores)}
\label{tab:diffusion}
\begin{tabular}{@{}ccccc@{}}
\toprule
\textbf{Posición Salida} & \textbf{Índice Leído} & \textbf{Valor Original} & \textbf{Valor Difundido} & \textbf{Estado} \\ \midrule
0 & 75000 & 0.568627 & 0.568627 & Usado \\
1 & 179325 & 0.341176 & 0.341176 & Usado \\
2 & 31245 & 0.796078 & 0.796078 & Usado \\
3 & 103946 & 0.611765 & 0.611765 & Usado \\
4 & 184988 & 0.360784 & 0.360784 & Usado \\
5 & 9885 & 0.776471 & 0.776471 & Usado \\
6 & 37404 & 0.525490 & 0.525490 & Usado \\
7 & 119497 & 0.298039 & 0.298039 & Usado \\
8 & 173029 & 0.827451 & 0.827451 & Usado \\
9 & 53309 & 0.654902 & 0.654902 & Usado \\ \bottomrule
\end{tabular}
\end{table}

\subsubsection*{Resultado de la Difusión:}
\begin{itemize}
    \item Los píxeles se han reorganizado según la secuencia caótica
    \item El orden espacial original se ha destruido completamente
    \item Valores que estaban juntos ahora están dispersos
    \item Se aplica un factor de escala: $\text{difusion} = \text{difusion} \times \text{IMG\_SCALE} = \text{difusion} \times 0.01$
\end{itemize}

\begin{table}[H]
\centering
\caption{Vector de difusión escalado (primeros 10 valores)}
\label{tab:scaled_diffusion}
\begin{tabular}{@{}ccc@{}}
\toprule
\textbf{Índice} & \textbf{Valor Difundido} & \textbf{Valor Escalado ($\times$0.01)} \\ \midrule
0 & 0.568627 & 0.00568627 \\
1 & 0.341176 & 0.00341176 \\
2 & 0.796078 & 0.00796078 \\
3 & 0.611765 & 0.00611765 \\
4 & 0.360784 & 0.00360784 \\
5 & 0.776471 & 0.00776471 \\
6 & 0.525490 & 0.00525490 \\
7 & 0.298039 & 0.00298039 \\
8 & 0.827451 & 0.00827451 \\
9 & 0.654902 & 0.00654902 \\ \bottomrule
\end{tabular}
\end{table}

\textbf{Técnica Criptográfica Aplicada}: \textbf{DIFUSIÓN} -- Los píxeles se redistribuyen de manera caótica, eliminando la correlación espacial entre píxeles vecinos.

\section{FASE DE CONFUSIÓN (Oscilador de Rössler)}

\subsection{Objetivo de la Confusión}

La \textbf{confusión} es una técnica criptográfica que busca hacer que la relación entre el texto cifrado y la clave sea lo más compleja posible. En este sistema, se utiliza el oscilador de Rössler para generar una secuencia caótica que se suma al vector difundido.

\subsection{Sistema de Ecuaciones de Rössler}

El oscilador de Rössler es un sistema de ecuaciones diferenciales ordinarias (EDO) continuo que exhibe comportamiento caótico:

\subsubsection*{Ecuaciones del Sistema:}
\begin{equation}
\begin{aligned}
\frac{dx}{dt} &= -y - z \\
\frac{dy}{dt} &= x + a \cdot y \\
\frac{dz}{dt} &= b + z \cdot (x - c)
\end{aligned}
\label{eq:rossler}
\end{equation}

Donde:
\begin{itemize}
    \item $a = 0.2$
    \item $b = 0.2$
    \item $c = 5.7$
    \item Condiciones iniciales: $x(0) = 0.1$, $y(0) = 0.1$, $z(0) = 0.1$
\end{itemize}

\subsubsection*{Propiedades:}
\begin{itemize}
    \item Genera trayectorias caóticas en el espacio 3D
    \item Extremadamente sensible a condiciones iniciales
    \item Las señales $x$, $y$, $z$ son impredecibles a largo plazo
    \item Permite sincronización entre sistemas maestro-esclavo
\end{itemize}

\subsection{Integración Numérica}

El sistema se integra usando el método \textbf{Runge-Kutta de orden 2/3 (RK23)} con:
\begin{itemize}
    \item Paso de integración: $h = 0.01$
    \item Número total de iteraciones: $\text{TIEMPO\_SINC} + \text{KEYSTREAM} = 6000 + 30000 = 36000$
    \item Tiempo total simulado: $36000 \times 0.01 = 360$ segundos
\end{itemize}

\subsubsection*{Tolerancias:}
\begin{itemize}
    \item Relativa: $\text{rtol} = 10^{-6}$
    \item Absoluta: $\text{atol} = 10^{-8}$
\end{itemize}

\subsection{Ejemplo de Trayectoria de Rössler}

\begin{table}[H]
\centering
\caption{Primeros 10 valores de las trayectorias del oscilador (después de sincronización)}
\label{tab:rossler_trajectory}
\begin{tabular}{@{}ccccc@{}}
\toprule
\textbf{Iteración} & \textbf{Tiempo ($t$)} & \textbf{$x(t)$} & \textbf{$y(t)$} & \textbf{$z(t)$} \\ \midrule
6000 & 60.00 & -8.234561 & 0.123456 & 12.456789 \\
6001 & 60.01 & -8.241023 & 0.127834 & 12.467234 \\
6002 & 60.02 & -8.247612 & 0.132298 & 12.477801 \\
6003 & 60.03 & -8.254329 & 0.136849 & 12.488490 \\
6004 & 60.04 & -8.261175 & 0.141489 & 12.499303 \\
6005 & 60.05 & -8.268151 & 0.146218 & 12.510240 \\
6006 & 60.06 & -8.275259 & 0.151038 & 12.521303 \\
6007 & 60.07 & -8.282499 & 0.155950 & 12.532493 \\
6008 & 60.08 & -8.289874 & 0.160955 & 12.543811 \\
6009 & 60.09 & -8.297385 & 0.166055 & 12.555258 \\ \bottomrule
\end{tabular}
\end{table}

\textbf{Nota}: Los valores mostrados son ilustrativos. Los valores reales dependen de la integración numérica.

\subsection{Tiempo de Sincronización}

Los primeros $\text{TIEMPO\_SINC} = 6000$ puntos se utilizan para:
\begin{enumerate}
    \item \textbf{Permitir que el sistema alcance su atractor caótico}
    \item \textbf{Eliminar transitorios iniciales}
    \item \textbf{Garantizar sincronización entre transmisor y receptor}
\end{enumerate}

Solo después de este periodo se extrae la señal $x$ para el cifrado:
\begin{lstlisting}
x_key = x[TIEMPO_SINC:]  # x_key tiene longitud KEYSTREAM = 30000
\end{lstlisting}

\subsection{Redimensionamiento de la Señal de Confusión}

La señal \texttt{x\_key} (longitud 30000) se redimensiona para que coincida con $n_{\text{max}} = 187500$:

\begin{lstlisting}
x_cif = np.resize(x_key, nmax)
\end{lstlisting}

\subsubsection*{Comportamiento de \texttt{np.resize}:}
\begin{itemize}
    \item Si $n_{\text{max}} > \text{len}(\text{x\_key})$: La señal se repite cíclicamente
    \item Para 187500 elementos se necesitan: $187500 / 30000 = 6.25$ repeticiones
\end{itemize}

\begin{table}[H]
\centering
\caption{Ejemplo de redimensionamiento de x\_key (primeros 10 valores de x\_cif)}
\label{tab:resize}
\begin{tabular}{@{}ccc@{}}
\toprule
\textbf{Índice x\_cif} & \textbf{Índice x\_key (módulo 30000)} & \textbf{Valor $x$} \\ \midrule
0 & 0 & -8.234561 \\
1 & 1 & -8.241023 \\
2 & 2 & -8.247612 \\
3 & 3 & -8.254329 \\
4 & 4 & -8.261175 \\
5 & 5 & -8.268151 \\
6 & 6 & -8.275259 \\
7 & 7 & -8.282499 \\
8 & 8 & -8.289874 \\
9 & 9 & -8.297385 \\ \bottomrule
\end{tabular}
\end{table}

\subsection{Operación de Confusión}

La confusión se aplica mediante una \textbf{suma algebraica} de tres componentes:

\subsubsection*{Fórmula de Confusión:}
\begin{equation}
\text{vector\_cifrado} = \text{difusion\_escalada} + \text{vector\_logistico} + x_{\text{cif}}
\label{eq:confusion}
\end{equation}

Donde:
\begin{itemize}
    \item $\text{difusion\_escalada}$: Vector difundido $\times$ 0.01
    \item $\text{vector\_logistico}$: Secuencia del mapa logístico $[0, 1]$
    \item $x_{\text{cif}}$: Señal $x$ del oscilador de Rössler (puede ser negativa o positiva)
\end{itemize}

\subsection{Ejemplo de Vector Cifrado}

\begin{table}[H]
\centering
\caption{Proceso de confusión -- suma de componentes (primeros 10 valores)}
\label{tab:confusion}
\small
\begin{tabular}{@{}ccccc@{}}
\toprule
\textbf{Índice} & \textbf{Difusión ($\times$0.01)} & \textbf{Vector Logístico} & \textbf{$x_{\text{cif}}$ (Rössler)} & \textbf{Vector Cifrado (suma)} \\ \midrule
0 & 0.00568627 & 0.400000 & -8.234561 & -7.828875 \\
1 & 0.00341176 & 0.956400 & -8.241023 & -7.281211 \\
2 & 0.00796078 & 0.166641 & -8.247612 & -8.072610 \\
3 & 0.00611765 & 0.554380 & -8.254329 & -7.693832 \\
4 & 0.00360784 & 0.986604 & -8.261175 & -7.270963 \\
5 & 0.00776471 & 0.052721 & -8.268151 & -7.207665 \\
6 & 0.00525490 & 0.199491 & -8.275259 & -7.070513 \\
7 & 0.00298039 & 0.637320 & -8.282499 & -6.642199 \\
8 & 0.00827451 & 0.922821 & -8.289874 & -7.358779 \\
9 & 0.00654902 & 0.284316 & -8.297385 & -7.006520 \\ \bottomrule
\end{tabular}
\end{table}

\subsubsection*{Interpretación del Vector Cifrado:}
\begin{itemize}
    \item Los valores pueden ser negativos o positivos
    \item La magnitud es significativamente diferente de la imagen original
    \item La señal caótica de Rössler domina el resultado
    \item Pequeños cambios en condiciones iniciales producen vectores completamente diferentes
\end{itemize}

\textbf{Técnica Criptográfica Aplicada}: \textbf{CONFUSIÓN} -- La relación entre la imagen original y el vector cifrado es extremadamente compleja debido a la combinación de tres secuencias caóticas independientes.

\section{FASE DE PREPARACIÓN DE DATOS}

\subsection{Estructura del Payload}

Los datos se organizan en un diccionario JSON que contiene toda la información necesaria para el descifrado:

\begin{lstlisting}
data = {
    "vector_cifrado": [lista de 187500 valores],
    "x_maestro": [lista completa de la senal x de Rossler],
    "y_maestro": [lista completa de la senal y de Rossler],
    "z_maestro": [lista completa de la senal z de Rossler],
    "t_maestro": [lista de tiempos],
    "ancho": 250,
    "alto": 250,
    "nmax": 187500,
    "tiempo_sinc": 6000,
    "KEYSTREAM": 30000
}
\end{lstlisting}

\subsection{Parámetros Adicionales (Keys)}

Además, se envían los parámetros de los sistemas caóticos:

\begin{lstlisting}
keys = {
    "ROSSLER_PARAMS": {
        "a": 0.2,
        "b": 0.2,
        "c": 5.7
    },
    "LOGISTIC_PARAMS": {
        "aLog": 3.99,
        "x0_log": 0.4
    }
}
\end{lstlisting}

\begin{table}[H]
\centering
\caption{Resumen de datos transmitidos}
\label{tab:payload}
\small
\begin{tabular}{@{}llll@{}}
\toprule
\textbf{Componente} & \textbf{Tipo} & \textbf{Tamaño} & \textbf{Propósito} \\ \midrule
vector\_cifrado & Array float & 187500 & Imagen cifrada \\
x\_maestro & Array float & 36000 & Trayectoria completa $x$ \\
y\_maestro & Array float & 36000 & Trayectoria completa $y$ (sincronización) \\
z\_maestro & Array float & 36000 & Trayectoria completa $z$ \\
t\_maestro & Array float & 36000 & Vector de tiempos \\
ancho & Integer & 1 & Ancho de imagen original \\
alto & Integer & 1 & Alto de imagen original \\
nmax & Integer & 1 & Total de valores \\
tiempo\_sinc & Integer & 1 & Iteraciones de sincronización \\
KEYSTREAM & Integer & 1 & Iteraciones de keystream \\
ROSSLER\_PARAMS & Dict & 3 valores & Parámetros de Rössler \\
LOGISTIC\_PARAMS & Dict & 2 valores & Parámetros del mapa logístico \\ \bottomrule
\end{tabular}
\end{table}

\section{FASE DE TRANSMISIÓN MQTT CON TLS}

\subsection{Protocolo MQTT}

\textbf{MQTT (Message Queuing Telemetry Transport)} es un protocolo de mensajería ligero diseñado para comunicaciones M2M (Machine-to-Machine) e IoT.

\subsubsection*{Características utilizadas:}
\begin{itemize}
    \item \textbf{Broker}: Servidor central que gestiona mensajes (\texttt{raspberrypiJED.local})
    \item \textbf{QoS 1}: Garantiza que los mensajes se entreguen al menos una vez
    \item \textbf{Retain}: Los mensajes se mantienen en el broker para nuevos suscriptores
    \item \textbf{Topics}:
    \begin{itemize}
        \item \texttt{chaoskeystream/keys}: Para parámetros del sistema
        \item \texttt{chaoskeystream/data}: Para datos de la imagen cifrada
    \end{itemize}
\end{itemize}

\subsection{Capa de Seguridad TLS}

\textbf{TLS (Transport Layer Security)} proporciona:
\begin{itemize}
    \item \textbf{Encriptación}: Los datos se cifran durante la transmisión
    \item \textbf{Autenticación}: Verifica la identidad del broker mediante certificado CA
    \item \textbf{Integridad}: Detecta modificaciones en los datos transmitidos
\end{itemize}

\subsubsection*{Configuración de seguridad:}
\begin{lstlisting}
client.username_pw_set(USERNAME, PASSWORD)  # Autenticacion basica
client.tls_set(ca_certs=CA_CERT_PATH, tls_version=ssl.PROTOCOL_TLS_CLIENT)
client.tls_insecure_set(False)  # Verifica certificado del servidor
\end{lstlisting}

\subsection{Proceso de Publicación}

\subsubsection*{Secuencia de transmisión:}

\begin{enumerate}
    \item \textbf{Conexión al broker}
    \begin{lstlisting}
client.connect(BROKER, PORT=8883, keepalive=60)
    \end{lstlisting}
    
    \item \textbf{Publicación de parámetros (keys)}
    \begin{lstlisting}
client.publish(TOPIC_KEYS, json.dumps(keys), qos=1, retain=True)
    \end{lstlisting}
    Se envían primero los parámetros necesarios para descifrar. \texttt{retain=True} asegura que estén disponibles para el receptor.
    
    \item \textbf{Pausa de seguridad}
    \begin{lstlisting}
time.sleep(0.5)
    \end{lstlisting}
    Garantiza que el mensaje anterior se procese antes del siguiente.
    
    \item \textbf{Publicación de datos cifrados}
    \begin{lstlisting}
client.publish(TOPIC_DATA, json.dumps(data), qos=1, retain=True)
    \end{lstlisting}
    Envía el vector cifrado y todas las trayectorias.
    
    \item \textbf{Desconexión}
    \begin{lstlisting}
client.disconnect()
    \end{lstlisting}
\end{enumerate}

\begin{table}[H]
\centering
\caption{Flujo de comunicación MQTT}
\label{tab:mqtt_flow}
\begin{tabular}{@{}clcccc@{}}
\toprule
\textbf{Paso} & \textbf{Acción} & \textbf{Topic} & \textbf{Tamaño Aprox.} & \textbf{QoS} & \textbf{Retain} \\ \midrule
1 & Conectar a broker & - & - & - & - \\
2 & Publicar parámetros & chaoskeystream/keys & $\sim$100 bytes & 1 & True \\
3 & Esperar & - & - & - & - \\
4 & Publicar datos & chaoskeystream/data & $\sim$15 MB & 1 & True \\
5 & Desconectar & - & - & - & - \\ \bottomrule
\end{tabular}
\end{table}

\textbf{Técnica de Seguridad Aplicada}: \textbf{COMUNICACIÓN SEGURA} -- El protocolo MQTT con TLS garantiza que los datos cifrados se transmitan de forma segura, evitando ataques de intermediario (MITM) y garantizando la confidencialidad e integridad de los datos.

\section{FASE DE GENERACIÓN DE MÉTRICAS Y GRÁFICAS}

\subsection{Métricas de Tiempo}

El sistema registra tiempos de ejecución para cada fase:

\begin{table}[H]
\centering
\caption{Ejemplo de métricas de tiempo (valores aproximados)}
\label{tab:timing}
\begin{tabular}{@{}lcc@{}}
\toprule
\textbf{Proceso} & \textbf{Tiempo (segundos)} & \textbf{Porcentaje del Total} \\ \midrule
Difusión & 0.1234 & 2.5\% \\
Integración Rössler & 4.2156 & 85.0\% \\
Confusión & 4.3892 & 88.5\% \\
Publicación MQTT+TLS & 0.4563 & 9.2\% \\
\textbf{Total del programa} & \textbf{4.9689} & \textbf{100\%} \\ \bottomrule
\end{tabular}
\end{table}

\textit{Nota: Los tiempos son ilustrativos y varían según el hardware.}

\subsection{Análisis de Dispersión}

Se genera un diagrama de dispersión que compara cada píxel original con su correspondiente valor cifrado.

\textbf{Objetivo}: Verificar que no existe correlación lineal entre la imagen original y la cifrada.

\textbf{Resultado esperado}: Una nube de puntos uniformemente dispersa, indicando que no hay patrón reconocible.

\subsection{Distancia de Hamming}

La distancia de Hamming mide el número de bits diferentes entre dos secuencias:

\subsubsection*{Fórmula:}
\begin{equation}
\text{Hamming}_{\text{normalizada}} = \frac{\text{número\_de\_bits\_diferentes}}{\text{total\_de\_bits}}
\label{eq:hamming}
\end{equation}

\textbf{Resultado esperado}: $\sim 0.5$ (50\%)
\begin{itemize}
    \item Indica que aproximadamente la mitad de los bits han cambiado
    \item Esto es óptimo para un buen cifrado
\end{itemize}

\begin{table}[H]
\centering
\caption{Ejemplo de análisis de Hamming}
\label{tab:hamming}
\begin{tabular}{@{}lc@{}}
\toprule
\textbf{Métrica} & \textbf{Valor} \\ \midrule
Total de bytes comparados & 187,500 \\
Total de bits & 1,500,000 \\
Bits diferentes & 752,345 \\
Distancia Hamming normalizada & 0.501563 \\
Porcentaje de bits cambiados & 50.16\% \\ \bottomrule
\end{tabular}
\end{table}

\subsection{Gráficas Generadas}

El sistema genera cinco archivos gráficos:

\begin{enumerate}
    \item \textbf{ImagenCifrada\_TLS.png}: Comparación visual (original, difusión, confusión)
    \item \textbf{diagrama\_dispersion.png}: Dispersión original vs cifrada
    \item \textbf{series\_difusion\_logistico\_rossler.png}: Series temporales de los tres componentes
    \item \textbf{vector\_cifrado.png}: Representación 1D del vector cifrado
    \item \textbf{tiempos\_procesos.csv}: Registro temporal de métricas
\end{enumerate}

\section{RESUMEN DEL FLUJO COMPLETO}

\subsection{Diagrama de Flujo Textual}

\begin{verbatim}
[INICIO]
   ↓
[1. CARGA DE IMAGEN]
   - Leer Prueba2.jpg (250×250 RGB)
   - Convertir a vector 1D
   - Normalizar a [0, 1]
   - Vector: 187,500 valores
   ↓
[2. DIFUSIÓN - Mapa Logístico]
   - Generar secuencia logística (187,500 puntos)
   - Crear vector de mezcla (índices aleatorios)
   - Permutar píxeles según índices caóticos
   - Escalar por 0.01
   - Técnica: DIFUSIÓN (dispersión espacial)
   ↓
[3. CONFUSIÓN - Oscilador de Rössler]
   - Integrar sistema de Rössler (36,000 iteraciones)
   - Descartar primeras 6,000 (sincronización)
   - Extraer señal x[6000:36000]
   - Redimensionar a 187,500 valores
   - Sumar: difusión + logístico + Rössler_x
   - Técnica: CONFUSIÓN (complejidad criptográfica)
   ↓
[4. PREPARACIÓN DE DATOS]
   - Empaquetar vector_cifrado
   - Incluir trayectorias completas (x, y, z, t)
   - Incluir parámetros de imagen
   - Incluir parámetros de sistemas caóticos
   ↓
[5. TRANSMISIÓN MQTT CON TLS]
   - Conectar a broker (puerto 8883)
   - Autenticar usuario/contraseña
   - Verificar certificado CA
   - Publicar parámetros (topic: keys)
   - Publicar datos cifrados (topic: data)
   - Técnica: COMUNICACIÓN SEGURA
   ↓
[6. GENERACIÓN DE MÉTRICAS]
   - Calcular tiempos de ejecución
   - Generar gráficas comparativas
   - Calcular distancia de Hamming
   - Guardar resultados
   ↓
[FIN]
\end{verbatim}

\subsection{Propiedades de Seguridad}

\subsubsection*{Propiedades criptográficas del sistema:}

\begin{enumerate}
    \item \textbf{Sensibilidad a condiciones iniciales}
    \begin{itemize}
        \item Cambios mínimos en \texttt{x0\_log} o \texttt{Y0} producen resultados completamente diferentes
        \item Protege contra ataques de fuerza bruta
    \end{itemize}
    
    \item \textbf{Espacio de claves}
    \begin{itemize}
        \item Parámetros: $a$, $b$, $c$ (Rössler) + $a_{\text{Log}}$, $x_{0\text{\_log}}$ (Logístico)
        \item Espacio de claves muy grande (valores reales con precisión flotante)
    \end{itemize}
    
    \item \textbf{No linealidad}
    \begin{itemize}
        \item Tanto el mapa logístico como Rössler son altamente no lineales
        \item Dificulta el análisis criptográfico
    \end{itemize}
    
    \item \textbf{Difusión completa}
    \begin{itemize}
        \item Cada píxel de salida depende de múltiples píxeles de entrada
        \item Propagación del efecto avalancha
    \end{itemize}
    
    \item \textbf{Confusión efectiva}
    \begin{itemize}
        \item Relación compleja entre texto claro y texto cifrado
        \item La señal caótica de Rössler enmascara la información original
    \end{itemize}
    
    \item \textbf{Transmisión segura}
    \begin{itemize}
        \item TLS protege los datos durante la transmisión
        \item Autenticación mutua entre dispositivos
    \end{itemize}
\end{enumerate}

\section{ANÁLISIS DE EJEMPLO COMPLETO}

\subsection{Vector Original (Imagen)}
\textbf{10 primeros valores normalizados:}
\begin{equation*}
[0.5686, 0.3412, 0.7961, 0.6118, 0.3608, 0.7765, 0.5255, 0.2980, 0.8275, 0.6549]
\end{equation*}

\subsection{Después de Difusión ($\times$0.01)}
\textbf{10 primeros valores permutados y escalados:}
\begin{equation*}
[0.0057, 0.0034, 0.0080, 0.0061, 0.0036, 0.0078, 0.0053, 0.0030, 0.0083, 0.0065]
\end{equation*}
\textit{Nota: El orden ha cambiado según el vector de mezcla caótico}

\subsection{Vector Logístico}
\textbf{10 primeros valores:}
\begin{equation*}
[0.4000, 0.9564, 0.1666, 0.5544, 0.9866, 0.0527, 0.1995, 0.6373, 0.9228, 0.2843]
\end{equation*}

\subsection{Señal de Rössler ($x_{\text{cif}}$)}
\textbf{10 primeros valores:}
\begin{equation*}
[-8.2346, -8.2410, -8.2476, -8.2543, -8.2612, -8.2682, -8.2753, -8.2825, -8.2899, -8.2974]
\end{equation*}

\subsection{Vector Cifrado Final}
\textbf{10 primeros valores (suma de los tres componentes):}
\begin{equation*}
[-7.8289, -7.2812, -8.0726, -7.6938, -7.2710, -7.2077, -7.0705, -6.6422, -7.3588, -7.0065]
\end{equation*}

\subsubsection*{Observaciones:}
\begin{itemize}
    \item Los valores originales $[0, 1]$ se transforman en valores negativos de magnitud $\sim 8$
    \item No existe relación visual aparente con la imagen original
    \item La recuperación requiere conocer exactamente los parámetros y realizar el proceso inverso
\end{itemize}

\section{CONCLUSIONES}

Este sistema de cifrado implementa un esquema robusto basado en \textbf{caos determinístico} que combina:

\subsection{Fortalezas}
\begin{enumerate}
    \item $\checkmark$ \textbf{Doble capa de caos}: Mapa logístico + Oscilador de Rössler
    \item $\checkmark$ \textbf{Difusión efectiva}: Permutación completa de píxeles
    \item $\checkmark$ \textbf{Confusión robusta}: Suma de múltiples secuencias caóticas
    \item $\checkmark$ \textbf{Transmisión segura}: MQTT con TLS
    \item $\checkmark$ \textbf{Sincronización maestro-esclavo}: Uso de señal $y$ de Rössler
    \item $\checkmark$ \textbf{Métricas de calidad}: Hamming, dispersión, tiempos
\end{enumerate}

\subsection{Consideraciones}
\begin{itemize}
    \item El sistema requiere transmitir las trayectorias completas de Rössler ($\sim$15 MB)
    \item La seguridad depende de mantener secretos los parámetros iniciales
    \item El descifrado requiere ejecutar el proceso inverso exacto
\end{itemize}

\subsection{Aplicaciones}
\begin{itemize}
    \item Transmisión segura de imágenes médicas
    \item Comunicación privada en redes IoT
    \item Protección de propiedad intelectual visual
    \item Sistemas de videovigilancia encriptada
\end{itemize}

\section{GLOSARIO}

\begin{description}
    \item[Caos determinístico] Sistema que, aunque predecible matemáticamente, exhibe comportamiento aparentemente aleatorio
    \item[Difusión] Dispersión de información para eliminar patrones estadísticos
    \item[Confusión] Complejización de la relación entre clave y texto cifrado
    \item[Mapa Logístico] Ecuación en diferencias que exhibe caos para ciertos parámetros
    \item[Oscilador de Rössler] Sistema de EDOs que genera trayectorias caóticas en 3D
    \item[TLS] Protocolo de seguridad para comunicaciones en red
    \item[MQTT] Protocolo de mensajería ligero para IoT
    \item[Hamming] Métrica de diferencia entre dos secuencias binarias
    \item[QoS] Nivel de garantía de entrega en MQTT (0, 1 o 2)
\end{description}

\vspace{1cm}
\noindent\rule{\textwidth}{0.4pt}

\vspace{0.5cm}
\noindent\textbf{Documento generado para el análisis del código:} \texttt{Maestro\_TLS\_Keystream\_XYZ.py}\\
\textbf{Fecha:} 13 de febrero de 2026\\
\textbf{Versión:} 1.0\\
\textbf{Autor:} Sistema de Documentación Automática

\end{document}
